\chapter{Campaigns, Clocks, and Consequences}
\label{ch:campaigns-clocks}
\index{campaigns}\index{clocks}\index{consequences}
\index{travel!GM guidance}

In \textbf{Fate's Edge}, campaigns are not a sequence of disconnected adventures---they are \textbf{pressure systems}.  
Every choice the players make pushes on something: a faction, a threat, a relationship, a belief. Over time, those pressures accumulate, shift, and eventually \emph{break}.

Your job as GM is not to hide this pressure---it is to \textbf{shape it, reveal it, and let the players decide where it lands}.

This chapter presents the core tools for doing that work: \textbf{Clocks}, \textbf{Fronts}, \textbf{the Crown Spread}, and long-term campaign structures. Used well, they create momentum without railroading and stakes without arbitrary escalation.

% =========================================================
% MOVED to Chapter 1 (Philosophy & GM as Player) - Session Zero Worksheet
% =========================================================
% \clearpage
% \section*{Session Zero Worksheet (GM)}
% \label{sec:gm-session-zero}
% 
% \begin{tcolorbox}[colback=blue!5!white,colframe=blue!75!black,title=Before the Session]
% Use this worksheet to capture the table's decisions. Keep it visible for the campaign.
% \end{tcolorbox}
% 
% \subsection*{1. Safety \& Tone}
% \begin{tabular}{p{4cm}p{10cm}}
% \textbf{Lines (hard no):} & \hrulefill \\
% \textbf{Veils (fade to black):} & \hrulefill \\
% \textbf{Tone (circle):} & Cinematic / Folk Horror / Political / Grim Survival / Swashbuckling \\
% \end{tabular}
% 
% \subsection*{2. Crown Spread (draw 5 cards)}
% \begin{tabular}{p{4cm}p{10cm}}
% \textbf{Spade (Site):} & \hrulefill \\
% \textbf{Heart (Rival):} & \hrulefill \\
% \textbf{Club (Pressure):} & \hrulefill \\
% \textbf{Diamond (Leverage):} & \hrulefill \\
% \textbf{Wild (Hidden):} & \hrulefill \\
% \end{tabular}
% 
% \noindent \emph{Interpret suits by regional flavour. Write one sentence per card.}
% 
% \subsection*{3. Campaign Clocks (visible)}
% \begin{tabular}{p{6cm}p{8cm}}
% \textbf{Mandate [6]} (party's public legitimacy): & \\
% \textbf{Crisis [6]} (opposition's growing strength): & \\
% \textbf{Rival's Agenda [6]}: & (name of rival party or faction) \\
% \end{tabular}
% 
% \subsection*{4. Character Hooks (one per player)}
% \begin{tabular}{p{4cm}p{10cm}}
% Player: & \hrulefill \\
% Character goal: & \hrulefill \\
% Character fear / debt: & \hrulefill \\
% Which Crown element are they tied to? & \hrulefill \\
% \end{tabular}
% 
% \subsection*{5. Patron Interest (if any)}
% \begin{tabular}{p{6cm}p{8cm}}
% Patron(s) watching the party: & \\
% Obligation style: & Passive (long campaign) / Active (short arc) \\
% \end{tabular}
% 
% \subsection*{6. First Scene Framing}
% \noindent \emph{Write an in media res opening line that ties to the Crown Spread:}
% \begin{tcolorbox}[colback=gray!5!white, colframe=gray!70!black, sharp corners]
% ``\hrulefill''
% \end{tcolorbox}
% 
% \subsection*{7. House Rules \& Optional Systems}
% \begin{tabular}{p{6cm}p{8cm}}
% Deck of Consequences? & Y / N \\
% Threat Pool? & Y / N \\
% Popcorn initiative? & Y / N \\
% Psionics allowed? & Y / N \\
% Other: & \\
% \end{tabular}
% \clearpage

% =========================================================
% CROWN SPREAD (stays in Chapter 8)
% =========================================================
\section{Crown Spread Cheat Sheet}
\label{sec:crown-spread}
\index{Crown Spread}\index{campaign!setup}

The Crown Spread is a five-card (or six-card) oracular tool for seeding a campaign's major elements. Use it during Session Zero or at the start of a new arc. It creates a web of relationships, not a railroad.

\subsection{What You Need}
\begin{itemize}
  \item A standard 52-card deck (jokers optional).
  \item A table or whiteboard to record the spread.
  \item 15--20 minutes of quiet thinking.
\end{itemize}

\subsection{The Spread Positions}
Draw one card for each position. Lay them in a semicircle (the \textquotedbl Crown\textquotedbl).

\begin{center}
\begin{tabular}{cll}
\toprule
\textbf{Position} & \textbf{Card Drawn} & \textbf{What It Represents} \\
\midrule
Spade & (e.g., 5\spadesuit) & \textbf{The Crown Site} -- where destiny is decided. \\
Heart & (e.g., Q\heartsuit) & \textbf{The Crown Rival} -- who stands between the party and their goals. \\
Club & (e.g., 9\clubsuit) & \textbf{The Crown Pressure} -- the relentless force that prevents complacency. \\
Diamond & (e.g., A\diamondsuit) & \textbf{The Crown Leverage} -- the advantage that can turn the tide. \\
Wild (draw a 6th card) & (e.g., K\spadesuit) & \textbf{The Hidden Force} -- the unexpected element. \\
\bottomrule
\end{tabular}
\end{center}

\subsection{Interpreting the Cards}
Use the suit to determine the nature of the element, then interpret the rank to refine it.

\begin{tcolorbox}[colback=gray!5!white, colframe=black, title=Quick Interpretation Guide]
\begin{description}
\item[Spade (Site)] -- A place: ruin, bridge, forest, tower, cave, city district, crossroads.
\item[Heart (Rival)] -- A person with emotional ties: family, former ally, rival, spurned lover, grieving widow.
\item[Club (Pressure)] -- A resource or systemic threat: famine, plague, debt, war tax, blight, curse.
\item[Diamond (Leverage)] -- An object or secret: a lost letter, a relic, a true name, a hidden passage, a witness.
\item[Wild (Hidden)] -- Anything: a monster, a traitor, a natural disaster, a forgotten god, a second faction.
\end{description}

\medskip
\noindent \textbf{Rank guidance:}
\begin{itemize}
  \item \textbf{Ace -- 3:} Minor scale, personal stakes.
  \item \textbf{4 -- 6:} Moderate scale, affects a community.
  \item \textbf{7 -- 10:} Major scale, threatens a region.
  \item \textbf{Jack -- King:} Face cards suggest a named individual with personality and agenda.
\end{itemize}
\end{tcolorbox}

\subsection{Worked Example (Silkstrand Campaign)}
\begin{quote}
\textbf{Draws:}
\begin{itemize}
  \item Spade: 5\spadesuit → Tithing Bridge (a small, working-class bridge in the dye district).
  \item Heart: J\heartsuit → Mira Venn, the Guildmistress's daughter, who once threw a rival off that bridge.
  \item Club: 9\clubsuit → The Rust -- a magical blight poisoning the dye vats.
  \item Diamond: A\diamondsuit → The Weeping Oracle, an ancient witness who knows the truth about the new red dye.
  \item Wild: K\spadesuit → The Drowned Man, a cursed dyer who haunts the canals.
\end{itemize}
\end{quote}

\subsection{Connecting the Elements}
After the draw, ask yourself:
\begin{enumerate}
  \item How does the Rival relate to the Site? (Mira's first kill was on the Tithing Bridge.)
  \item How does the Pressure affect the Site or Rival? (The Rust is poisoning the canal below the bridge; Mira needs the Rust to continue to control the dye market.)
  \item How can the Leverage help against the Pressure or Rival? (The Oracle knows the Rust's source and the true cost of Mira's dye.)
  \item What does the Hidden Force want, and how does it complicate the others? (The Drowned Man kills polluters – he is a wildcard enemy and potential ally.)
\end{enumerate}
Draw a simple relationship map. You now have a campaign skeleton.

\subsection{Clocks for the Crown Spread}
Create one visible clock for each element that will take time to resolve. Use these defaults:

\begin{center}
\begin{tabular}{cl}
\toprule
\textbf{Element} & \textbf{Sample Clock} \\
\midrule
Rival's Agenda & \textit{Mira's Consolidation} [6] \\
Pressure & \textit{The Rust Spreads} [8] \\
Leverage (if it can be lost) & \textit{The Oracle's Safety} [6] \\
Hidden Force (if it escalates) & \textit{The Drowned Man's Attention} [4] \\
\bottomrule
\end{tabular}
\end{center}

Tick these clocks during play when players fail relevant rolls or ignore the element. When a clock fills, that element's situation worsens dramatically.

\subsection{Seasonal Evolution (Optional)}
For longer campaigns, evolve the Crown Spread through seasons:

\begin{description}
\item[Winter (Establishment)] – Introduce elements as rumours or distant threats.
\item[Spring (Growth)] – Elements become active; the Rival makes a move.
\item[Summer (Climax)] – Confrontation at the Crown Site.
\item[Autumn (Harvest)] – Consequences manifest; the party's choices become legacy.
\end{description}

\subsection{GM Quick Cues}
\begin{itemize}
  \item \textbf{Don't over-plan.} The spread is a seed, not a script. Let players discover connections through play.
  \item \textbf{Be flexible.} If a better idea emerges, replace an element – the cards are oracles, not chains.
  \item \textbf{Use the Wild card for surprise.} Even you don't have to know what it is until it appears.
  \item \textbf{Tie every character to at least one element.} During Session Zero, ask each player: ``How do you know about the Rival?'' or ``Why does the Pressure affect you personally?''
\end{itemize}

\subsection{Crown Spread at a Glance – Ritual Card}
For a quick reference, print the following card and keep it behind your GM screen:

\begin{tcolorbox}[colback=yellow!5!white, colframe=black, title=Crown Spread – One-Card Ritual]
\begin{verbatim}
1. Name your campaign theme.
2. Shuffle, draw 5 cards (6 for Wild).
3. Layout: Spade (Site) | Heart (Rival) | Club (Pressure) | Diamond (Leverage) | (Wild)
4. Interpret suit: Place, Person, Scarcity, Secret, Twist.
5. Connect: How does each relate to the others?
6. Create 4 clocks (4–8 segments).
7. Play to find out.
\end{verbatim}
\end{tcolorbox}

% =========================================================
% MOVED to Chapter 1 (Philosophy & GM as Player) - Framework, Not a Prescription
% =========================================================
% \section{A Framework, Not a Prescription}
% \label{sec:framework-mindset}
% \index{GM mindset!modularity}
% 
% Before you build your campaign, internalise this truth: \textbf{You own the game. The game does not own you.}
% 
% Fate's Edge provides a library of tools -- clocks, fronts, travel procedures, faction turns, the Crown Spread, and many more. You are not expected to use all of them. In fact, you are encouraged to use only what your table needs \emph{right now}.
% 
% \begin{tcolorbox}[colback=gray!5!white, colframe=gray!75!black, title={The 10\% Rule}]
%   As the Grey Wanderer says: \textit{``A single page, well-used, is worth more than a library of unopened tomes. Pick one thing from one expansion. Use it tonight.''}
% 
%   The same applies to this chapter. Start with:
%   \begin{itemize}
%     \item \textbf{Clocks} -- always useful.
%     \item \textbf{The Crown Spread} -- for session zero.
%     \item \textbf{One front} -- to give the world a pulse.
%   \end{itemize}
%   
%   Ignore faction turns, legacy bonds, and travel sub-systems until they become relevant. 
%   Add them when you feel ready, not because the book says they exist.
% 
%   \textbf{Discard what doesn't serve your table.} Your campaign will not collapse if you never use the ``Clock Relationships'' section. It will not break if you skip the faction turn. The rules are here to serve you -- not the other way around.
% \end{tcolorbox}

% =========================================================
% MOVED to Chapter 4 (Cognitive Load) - Clocks section
% =========================================================
% \section{Clocks: Visible Pressure, Meaningful Choice}
% \label{sec:clocks}
% \index{campaign clocks}
% 
% ... (entire section with all subsections commented out) ...

% =========================================================
% Factions and Rivals (stays in Chapter 8)
% =========================================================
\section{Factions and Rivals}
\label{sec:factions-rivals}

A faction is any organised group with goals that intersect the party's. A rival party is a special kind of faction -- a small group of NPC adventurers who mirror the PCs.

\subsection{Faction Sheet (Quick Template)}
\label{subsec:faction-template}

\begin{tcolorbox}[colback=white!97!gray, colframe=black!80!gray, title=Faction Template]
\begin{tabular}{p{4cm}p{10cm}}
\textbf{Name:} & \hrulefill \\
\textbf{Agenda Clock:} & [\hspace{1cm}] (what they want to achieve) \\
\textbf{Key NPCs:} & (2--3 names and roles) \\
\textbf{Resources:} & (what they can bring to bear) \\
\textbf{Ticking conditions:} & (when does their clock advance?) \\
\textbf{Visible tracks:} & (list clocks players can see) \\
\end{tabular}
\end{tcolorbox}

\subsection{Rival Parties}
\label{subsec:rival-parties}

A rival party is a faction of 3--5 NPCs with their own Patrons, assets, and a visible \textbf{Rival Agenda clock} (6--8 segments). Tick their clock when the PCs ignore them or the GM spends SB on their off-screen activities. When the clock fills, the rivals achieve a major milestone. Reveal their progress through rumours, discovered scenes, or direct confrontation.

For a full template and example (The Ashen Cord), see the GM's Appendix \ref{app:rival-parties}.

% =========================================================
% Faction Turn (stays in Chapter 8)
% =========================================================
\section{Faction Turn Cheat Sheet}
\label{sec:faction-turn}
\index{factions!turn}\index{clocks!faction}

Between sessions (or at the start of a session), you can run a quick \textbf{faction turn} to keep the world alive. A faction turn takes 5--10 minutes and answers: \textit{``What have the major powers done while the party was busy?''}

\subsection{When to Run a Faction Turn}
\begin{itemize}
  \item At the end of a session, after players declare their next destination.
  \item At the beginning of a session, as a ``news from the world'' recap.
  \item During a long downtime (several weeks or months).
\end{itemize}

\subsection{Procedure – 5 Minutes}
\begin{enumerate}
  \item \textbf{List active factions.} Keep no more than 3--5 that directly interact with the party.
  \item \textbf{For each faction, roll 1d6.} Use the table below to determine their progress.
  \item \textbf{Advance or retreat their Agenda Clock.} Each faction has a visible clock (6--10 segments) representing their goal.
  \item \textbf{Note one concrete change.} Write one sentence per faction (e.g., ``The Scarlets seized the river docks.'').
  \item \textbf{Read the changes aloud} at the start of the next session.
\end{enumerate}

\begin{tcolorbox}[colback=blue!5!white, colframe=blue!75!black, title=Faction Progress Roll (d6)]
\begin{center}
\begin{tabular}{cll}
\toprule
\textbf{Roll} & \textbf{Progress} & \textbf{Agenda Clock Change} \\
\midrule
1–2 & Setback & –1 segment (min 0) \\
3–4 & Stalled & No change \\
5–6 & Advance & +1 segment \\
\bottomrule
\end{tabular}
\end{center}
\end{tcolorbox}

\subsection{Optional: Faction Conflict Roll}
If two factions are directly opposed (e.g., the Scarlets and the Grey Circle), roll opposed d6. Higher roll wins; lower roll loses 1 segment. Ties mean a stalemate – both lose 1 segment.

\subsection{Example Faction Turn}
\begin{quote}
\textbf{Factions:}
\begin{itemize}
  \item Mira's Scarlets – Agenda Clock: \textit{Control the Dye Trade} [8] (currently 3/8)
  \item Church of the Everflame – Agenda Clock: \textit{Suppress Heresy} [6] (currently 2/6)
  \item The Drowned Man (not a faction, but a wildcard) – no clock
\end{itemize}

\noindent \textbf{Rolls:}
\begin{itemize}
  \item Scarlets: 5 → Advance → now 4/8.
  \item Everflame: 2 → Setback → now 1/6.
\end{itemize}

\noindent \textbf{One-sentence changes:}
\begin{itemize}
  \item \textit{Scarlets:} ``Mira's agents have bought the loyalty of the harbour master.''
  \item \textit{Everflame:} ``An inquisitor was found murdered; the Church retreats to its cathedral.''
  \item \textit{Drowned Man:} (GM fiat) ``The Drowned Man dragged two Scarlets into the canal. Rumours spread.''
\end{itemize}

\noindent \textbf{Next session opener:} ``You wake to news that the harbour master has doubled tariffs – and that a body was pulled from the water wearing a Scarlets' token.''
\end{quote}

\subsection{Using Faction Clocks in Play}
\begin{itemize}
  \item \textbf{When a clock fills}, the faction achieves its goal. The party must react (or suffer the consequences).
  \item Players can \textbf{affect faction clocks} during adventures: a successful sabotage mission might tick a clock backwards; a public failure might tick it forward.
  \item Keep faction clocks \textbf{visible} (on an index card or whiteboard). Players should see the pressure building.
\end{itemize}

\subsection{Lightweight Faction Standing (Optional)}
For factions that are not pursuing a single agenda but merely have an attitude toward the party, track a simple \textbf{-3 to +3 scale}:
\begin{itemize}
  \item \textbf{-3:} Enemy – actively works against the party.
  \item \textbf{0:} Neutral – indifferent.
  \item \textbf{+3:} Ally – will sacrifice for the party.
\end{itemize}
Shift standing by ±1 after major interactions (saving a faction's leader, betraying them, etc.). No dice needed – just fiction.

\subsection{GM Quick Cues}
\begin{itemize}
  \item \textbf{Don't overdo it.} One faction turn per 1--2 sessions is plenty.
  \item \textbf{Tie changes to player actions.} If the party ignored a threat, tick that faction's clock for free.
  \item \textbf{Use the turn to generate rumours.} ``The Grey Circle is digging a new well – some say they found bones.''
\end{itemize}


% =========================================================
% PART 2: SEEDING THE CAMPAIGN (CROWN SPREAD ALREADY ABOVE)
% =========================================================
% (The Crown Spread sections appear earlier in this chapter; they are not duplicated here.)


% =========================================================
% PART 3: RUNNING THE CAMPAIGN
% =========================================================
\section{Running Travel as a Medieval Trek (Abridged)}
\label{sec:gm-travel}
\index{travel!GM techniques}

Travel in Fate's Edge is not a line on a map. It is a \textbf{medieval journey} -- uncertain, weather-beaten, dependent on local guides, hospitality, and the kindness (or cruelty) of strangers. For full procedures (weather tables, road condition dice, regional decks), see the \textit{Travel Guide} and \textit{Player's Guide} Chapter 10. This section gives you the GM principles.

\subsection{Lead with Weather, Not Distance}
Open each leg with a weather description, then set Position based on that weather (e.g., Desperate for navigation in fog, Controlled for stealth in rain). The weather \emph{is} the terrain.

\subsection{Maps Are Lies; Guides Are Truth}
Require a \textbf{Guide} role for each leg (PC or hired NPC). The Guide's roll (Wits + Survival or Lore) determines if they find the safe route, the short route, or the one that leads to a cultist ambush.

\subsection{Hospitality as a Resource}
When the party seeks shelter, draw a \textbf{Heart} from the regional deck -- that's the host. Their attitude affects what they offer. Spending a \textbf{String} (e.g., ``I helped this village last year'') can improve the outcome.

\subsection{Road Conditions Change Daily}
At the start of each leg, choose or roll a Road Condition (see the \textit{Travel Guide} for the full d6 table). Apply its effect for the entire leg.

\subsection{The Travel Clock as Mood Tracker}
Do not just tick the clock on successes. Use it to track the \emph{feeling} of the journey (Optimism $\rightarrow$ Fatigue $\rightarrow$ Desperation $\rightarrow$ The Edge). When the clock fills, describe the arrival -- mud-caked and grateful.

\subsection{Genre Shifts Through Borders}
When the party crosses from one region to another, the dominant fantasy subgenre shifts. Use the travel deck to introduce the new region's flavour immediately (e.g., from Ecktoria's political intrigue to the Mistlands' folk horror: the road grows silent, fog rolls in, a bell tolls without wind).

\subsection{Travelling Without a Map}
Fate's Edge does not use hex crawls. Travel is abstracted into \textbf{legs}, each with a Travel Clock (4, 6, 8, or 10 segments). You need only know how many \textbf{meaningful scenes} stand between two locations -- not the miles.

\section{Pacing and Arc Structure}
\label{sec:pacing-arcs}

A well-paced campaign follows a natural rhythm. Use these guidelines for session and arc structure.

\subsection{Session Energy Management}
Aim for three ``beats'' per session:
\begin{enumerate}
  \item \textbf{Opening hook} (5--10 minutes): Re-establish stakes, remind of clocks.
  \item \textbf{Rising action} (60--90 minutes): 2--3 scenes with escalating tension.
  \item \textbf{Climax or cliffhanger} (20--30 minutes): A major roll, a clock filling, or a hard choice.
\end{enumerate}

\subsection{Arc Length by Tier}
\begin{itemize}
  \item \textbf{Tier I--II (0--90 XP)}: 3--5 session arcs. Local scope, personal stakes.
  \item \textbf{Tier III (91--150 XP)}: 6--10 session arcs. Regional influence, branching paths.
  \item \textbf{Tiers IV--V (151+ XP)}: 15+ session arcs. Continental scope, world-shaping consequences.
\end{itemize}

\subsection{Between Sessions: The GM's Sacred Trust}
\label{sec:between-sessions}
\index{GM responsibilities}

Between sessions, update clocks, track complication debt, and prepare scenes. Use the \textit{Between-Session Checklist} in the GM's Appendix. Remember: your preparation is the quiet magic that transforms random encounters into meaningful episodes.

% =========================================================
% PART 4: THE FINALE AND BEYOND
% =========================================================
\section{Calling or Forcing the Crown}
\label{sec:calling-crown}
\index{Crown}\index{Finale}

The campaign reaches its crescendo when one of two thresholds is met---the moment when accumulated influence and mounting pressure collide in a final reckoning.

\begin{tcolorbox}[colback=white!97!gray, colframe=black!80!gray, title=Finale Triggers and Conditions]
\begin{tabularx}{\linewidth}{lX}
\toprule
\textbf{Finale Type} & \textbf{Conditions and Narrative Implications} \\
\midrule
Player-Called & Mandate $\geq$ 6 and Crisis $\leq$ 3 -- the party has earned the right to choose their moment of triumph. \\
Forced Finale & Crisis $\geq$ 6 regardless of Mandate -- the world forces a confrontation that can no longer be avoided. \\
Balanced Finale & Both dials at 4--5 -- a tense equilibrium where victory and defeat hang in perfect balance. \\
\bottomrule
\end{tabularx}
\end{tcolorbox}

\subsection{The Finale Procedure}
\label{subsec:finale-procedure}

When the Crown is called, run the three-beat finale:
\begin{enumerate}
  \item \textbf{Reckoning}: Defend or sanctify the record of accomplishments. Draw upon the Rival's established motives. Place the Pressure rail that will drive the scene forward.
  \item \textbf{Crossing}: Stage the kinetic rail (Escape/Hunt/Hazard) that threatens to end the scene prematurely if not managed carefully.
  \item \textbf{Coronation}: Use the Diamond Leverage to sign, seal, or swear the oath that cements the campaign's legacy.
\end{enumerate}

\subsection{Twist Collision (Finale Clause)}
\label{subsec:twist-collision}

Exactly once, when the Rival's Spade Twist contradicts their Club Belief, the table chooses:
\begin{itemize}
  \item GM gains +1 SB to complicate matters, or
  \item Players reduce two ticks total across the active rails, gaining breathing room.
\end{itemize}

\begin{tcolorbox}[colback=blue!5!white, colframe=blue!75!black, title={Twist Collision Example}]
  \textit{The Rival's Spade Twist:} ``The bell tower collapses, killing the rival's lieutenant.''  
  \textit{The Rival's Club Belief:} ``No more bloodshed -- I will end this peacefully.''  
  \textbf{Contradiction:} The rival now has no reason to hold back.  
  \textbf{Table chooses:} 
  \begin{itemize}
    \item GM gains +1 SB (the rival blames the PCs for the death), or
    \item Players reduce two ticks on the Collapse clock (they pull the rival from the rubble, saving a life).
  \end{itemize}
\end{tcolorbox}

\section{Legacy Conversion: Epilogue}
\label{sec:legacy-conversion}
\index{Legacy Conversion}\index{epilogue}

After the Finale, each PC draws 2 cards and answers epilogue prompts by suit. Then convert campaign elements into lasting legacy:
\begin{itemize}
  \item \textbf{Major Asset $\rightarrow$ Institution} (12 XP): A safehouse becomes a school, a spy ring becomes an intelligence service.
  \item \textbf{Seasonal Endorsement $\rightarrow$ Doctrine Rider} (4 XP): Temporary support becomes permanent policy.
  \item \textbf{Follower (Cap 3+) $\rightarrow$ Stationed NPC} (0 XP): Loyal companions become custodians of the new order.
  \item \textbf{Rival $\rightarrow$ Fixture}: Surviving adversaries become recurring elements of the setting's fabric.
\end{itemize}

\subsection{The Legacy Bond}
\label{subsec:legacy-bond}

When a PC retires or dies, their successor (a follower or a new character) inherits a \textbf{Legacy Bond}. This bond represents the predecessor's enduring influence. Once per session, the player may invoke the Legacy Bond to gain \textbf{1 Boon} when acting to uphold the predecessor's values or complete their unfinished business.

\textbf{Sample Legacy Bonds:}
\begin{itemize}
  \item ``I carry Kael's sword. I will not let it rust in peace.''
  \item ``Grandmother's garden still grows. I protect it as she did.''
  \item ``The oath we swore -- I will see it fulfilled.''
\end{itemize}

\subsection{Character Arc Management}
\label{subsec:character-arcs}

Help players develop meaningful character growth through the legacy process:
\begin{enumerate}
  \item \textbf{Establishment}: Define the character's current state and potential conflicts.
  \item \textbf{Development}: Create opportunities for growth and choice.
  \item \textbf{Crisis}: Present challenges that test the character's core beliefs.
  \item \textbf{Resolution}: Allow meaningful transformation based on choices.
\end{enumerate}

\begin{tcolorbox}[colback=green!5!white, colframe=green!75!black, title={Arcs Generate Obligation and Patron Interest}]
  \index{character arcs!obligation}\index{Patrons!in character arcs}
  A character's personal arc should not exist in a vacuum. Tie it directly to the game's resource systems:
  \begin{itemize}
    \item \textbf{Internal arcs} (fear, doubt, guilt) are perfect sources of \textit{Corruption} or \textit{Obligation}. A character struggling with guilt might be tempted by Malachai. A character seeking redemption might attract the attention of the Gallow's Bell.
    \item \textbf{External arcs} (revenge, ambition, love) create \textit{Debt Clocks} with terrestrial patrons or generate \textit{Crisis} when they interfere with faction plans.
    \item \textbf{Whenever a player advances their arc, mark +1 Obligation to a relevant Patron or tick a Patron's Interest clock.} The arc becomes a source of complications, not just backstory.
  \end{itemize}
  \textit{Example:} Sven wants to avenge his murdered brother. The GM notes that his brother was killed by a cult of Malachai. Each time Sven pursues a clue or confronts a cultist, the GM ticks \emph{Malachai's Interest [6]}. When the clock fills, Malachai appears in a dream and offers Sven the power to find the killer -- for a price.
\end{tcolorbox}

% =========================================================
% PART 5: QUICK START & EXAMPLES
% =========================================================
\section{Quick Start Campaigns}
\label{sec:quick-start-campaigns}

\subsection{The Sunstone Gambit (Tier I--II, 4--6 sessions)}
A medium-difficulty introduction to Fate's Edge set in Acasia.  
\textbf{Premise:} A scholar hires the party to retrieve a sunstone from a ruined tower before a rival adventuring company does.  
\textbf{Key NPCs:} Scholar Mira (quest-giver), Captain Voss (rival leader), the Tower's Guardian (ancient spirit).  
\textbf{Clocks:} \emph{Rival Progress [6]}, \emph{Tower Collapse [8]}, \emph{Spirit Wakening [4]}.  
\textbf{Session breakdown:} Commission $\rightarrow$ Travel (mountain leg) $\rightarrow$ Exploration (tower rooms) $\rightarrow$ Confrontation (rival team) $\rightarrow$ Climax (the Guardian) $\rightarrow$ Aftermath.

\subsection{The Dying Sun (Tier III--IV, 6--8 sessions)}
A regional crisis campaign. The sun-spirit is fading; a cult of eternal night seeks to extinguish it permanently.  
\textbf{Premise:} The party must rally factions, recover a lost ritual, and confront the cult before the winter solstice.  
\textbf{Key NPCs:} High Priestess of the Sun, Cult Shadow-bound, Reluctant Ykrul Chieftain.  
\textbf{Clocks:} \emph{Sun's Fading [10]} (ticks each session), \emph{Cult Influence [6]}, \emph{Alliance of the Dawn [8]}.  
\textbf{Acts:} Investigate the cult $\rightarrow$ Gather allies $\rightarrow$ Raid the cult's lair $\rightarrow$ Final confrontation.

\subsection{Tier-Appropriate Design}
\label{subsec:tier-design}
\begin{itemize}
  \item \textbf{Tier I--II (0--90 XP)}: Local scope, personal stakes, 3--5 session arcs.
  \item \textbf{Tier III (91--150 XP)}: Regional influence, branching arcs, 6--10 sessions.
  \item \textbf{Tiers IV--V (151+ XP)}: Continental scope, world-shaping consequences, multi-arc campaigns (15+ sessions).
\end{itemize}
For deeper guidance on high-tier campaigns, see Chapter \ref{ch:gm-epic} (\textit{Tier IV and V: Running Epic Adventures}).


% =========================================================
% CAMPAIGN WORKSHEET (One Page)
% =========================================================
\section{Campaign Worksheet}
\label{sec:worksheet}

Use this one-page worksheet to capture your campaign's key elements.

\begin{tcolorbox}[colback=white!97!gray, colframe=black!80!gray, title=Campaign Worksheet]
\begin{tabular}{p{4cm}p{10cm}}
\textbf{Campaign Name:} & \hrulefill \\
\textbf{Starting Tier:} & \hrulefill \quad \textbf{Estimated Sessions:} \hrulefill \\
\end{tabular}
\vspace{2mm}
\textbf{Crown Spread:} \\
\begin{tabular}{ll}
Spade (Site): & \hrulefill \\
Heart (Rival): & \hrulefill \\
Club (Pressure): & \hrulefill \\
Diamond (Leverage): & \hrulefill \\
Wild (Hidden): & \hrulefill \\
\end{tabular}
\vspace{2mm}
\textbf{Factions (name, agenda clock, key NPCs):} \\
\hrulefill \\
\hrulefill \\
\vspace{2mm}
\textbf{Fronts (name, agenda clock, supporting clocks, stakes):} \\
\hrulefill \\
\hrulefill \\
\vspace{2mm}
\textbf{Mandate:} 0/6 \quad \textbf{Crisis:} 0/6 \\
\vspace{2mm}
\textbf{Notes:} \hrulefill
\end{tcolorbox}


% =========================================================
% FINAL NARRATIVE REMINDER
% =========================================================
\section{Narrative First: The World Remembers}
\label{sec:narrative-first}
\index{narrative first}

Campaign design in Fate's Edge is not about railroading players along predetermined paths---it's about \textbf{responding to player choices} with consequences that accumulate like stones in a river, gradually shaping the flow of the narrative itself. Let the world shift in response to their actions. Let factions rise and fall based on their allegiances. Let the dice sing the song of a universe that reacts.

And when the Crown is finally crowned---when the last card is played and the final clock ticks to completion---let the echo of that moment be heard across the entire Amaranthine, a testament to stories well-lived and consequences fully earned.

Remember: Your preparation between sessions is the quiet magic that transforms random encounters into meaningful episodes and mechanical challenges into memorable stories. The investment in this sacred trust pays dividends in player engagement, narrative coherence, and the creation of campaigns that will be remembered long after the final dice have been rolled.